\section{What Federation Costs}
\label{sec:ch01-what-federation-costs}

\index{federation!cost of}%
The first cost is a contract every service signs and keeps signing. To act as a
subgraph, a service must extend its schema with the federation definitions,
resolve \code{Query.\_service}, provide a mechanism for resolving entity fields
through \code{Query.\_entities}, and apply \code{@link} to its \code{schema}
type to opt into Federation v2~\autocite{apollosubgraphspec}. None of those
four will mean anything to you yet, and they are not supposed to:
chapter~\ref{ch:first-subgraph} is where you implement them and see what each
one does. Count them for now. Four requirements, in every service,
permanently. A library implements them for you and this book uses one, which
changes who writes the code and not who owns the obligation.

The second is that composition becomes a build step that can fail on a change
you did not make. Two subgraphs that each compile can still refuse to compose,
because the schemas disagree about a type they share. The failure surfaces
after both teams have merged, it names a route through the graph rather than a
line in a file, and reading those messages is a skill.
Chapter~\ref{ch:satisfiability} is about the ones that are genuinely hard to
read, and appendix~\ref{app:comperrors} decodes the rest.

The third is that you now operate a router, and it is on the request path for
everything. Ernst's team at Booking.com reports \enquote{nearly 1,000 instances
of our Gateway running in production}~\autocite{ernst2022}. That is what the
tier becomes at their scale. Yours will be smaller, and it will still be a tier
you capacity-plan, monitor, roll out and get paged for, which you did not have
before.

The fourth is quieter and it is the one I would weigh most heavily. Running the
graph locally now means running several services, and answering \enquote{why is
this field null} means reading traces across all of them.
Chapter~\ref{ch:blame-routing} is entirely about the fact that when a federated
graph misbehaves, the team that owns the router gets asked first and usually
cannot answer.

The largest cost is not on that list. The coordination did not disappear when
you moved the merge into a program; it moved into the things the program
checks. Which service owns which type, which fields make up a key, which
subgraph is allowed to extend somebody else's entity: these are negotiated
between teams, the same teams as before, and the negotiation is now
load-bearing because the composer enforces the result.
Chapter~\ref{ch:ownership} gives that its own chapter, and it is there because
I think ownership is where federated graphs actually come apart, well before
anyone struggles with the composer.
